\subsection{Global Mean Thresholding}

The first dynamic method calculates one reference intensity from the
complete grayscale image. If the grayscale matrix is denoted by $Y$,
the reference is the mean intensity of all pixels:
\begin{align}
\mu &= \frac{1}{MN}\sum_{i=1}^{M}\sum_{j=1}^{N}Y(i,j)\label{eq:global_mean}
\end{align}

Every pixel is then compared with this single image-wide reference. The
corresponding binary decision is
\begin{align}
B(i,j)&=
\begin{cases}
255, & Y(i,j)\geq\mu,\\
0, & Y(i,j)<\mu.
\end{cases}\label{eq:global_mean_threshold}
\end{align}

This method is adaptive in the sense that the reference changes with the
input image, but the reference itself is global. It therefore does not
change from one location to another within the same image.
\vspace*{\baselineskip}
\subsection{Static Reference Thresholding}

For comparison, a fixed reference of $128$ is used in the static
thresholding method. A grayscale pixel is assigned to white when its
intensity is at least $128$ and to black otherwise:
\begin{align}
B(i,j)&=
\begin{cases}
255, & Y(i,j)\geq128,\\
0, & Y(i,j)<128.
\end{cases}\label{eq:static_threshold}
\end{align}

The purpose of this fixed reference is to provide a simple baseline
against which the dynamically calculated references can be compared.
\vspace*{\baselineskip}
\subsection{Gaussian-Weighted Local Reference}

A more localized reference is obtained by applying a Gaussian filter to
the grayscale intensity matrix. Instead of giving every neighbouring
pixel equal importance, the Gaussian function gives greater importance
to pixels close to the point being considered and progressively less
importance to pixels farther away.

The Gaussian weighting function is
\begin{align}
w(d)&=\exp\left(-\frac{d^2}{2\sigma^2}\right)\label{eq:gaussian_weight}
\end{align}

Here, $d$ is the distance from the point being processed and $\sigma$
controls the spatial spread of the Gaussian weighting. The corresponding
local reference can be understood as the weighted average
\begin{align}
\mu_G(i,j)&=
\frac{\displaystyle\sum_{(u,v)\in\mathcal{N}}
w(d_{uv})Y(u,v)}
{\displaystyle\sum_{(u,v)\in\mathcal{N}}w(d_{uv})}\label{eq:gaussian_local_mean}
\end{align}

Thus, the reference intensity is different at different locations in
the image and follows the local intensity distribution smoothly.

In the implementation, this operation is performed using
\texttt{gaussian\_filter}:
\newpage
\begin{lstlisting}[frame=single, basicstyle=\rmfamily]
local_mean = gaussian_filter(
    grayscale_matrix.astype(np.float32), sigma=sigma
)
\end{lstlisting}
\vspace*{\baselineskip}
\subsection{Importance of $\sigma$}

The value of $\sigma$ determines how localized the Gaussian reference
is. A smaller value concentrates the weighting around the current pixel,
so the reference responds strongly to nearby intensity variations. A
larger value spreads the weighting over a wider region, causing pixels
farther from the current point to contribute more significantly.

Consequently, $\sigma=10$ gives the localized reference used in this
implementation, whereas $\sigma=20$ gives the wider reference. As
$\sigma$ becomes sufficiently large, the Gaussian weights over the
relevant region become increasingly similar. The Gaussian-weighted
average therefore approaches the behaviour of an ordinary average over a
broad neighbourhood.
\vspace*{\baselineskip}
\subsection{Purpose of $C=3$}

After the Gaussian local reference is calculated, the implementation
uses the threshold
\begin{align}
T_G(i,j)&=\mu_G(i,j)-C\label{eq:gaussian_threshold}
\end{align}

The constant $C=3$ is introduced as a small offset to the Gaussian local
reference. Subtracting this value slightly lowers the local threshold,
allowing pixels that are marginally darker than the calculated local mean
to still be classified as white. In this implementation, $C=3$ is used
as a small empirical adjustment to fine-tune the sensitivity of the
resulting black-and-white output.
\vspace*{\baselineskip}
\subsection{Gaussian Local Thresholding}

The complete decision made by the Gaussian method is therefore
\begin{align}
B(i,j)&=
\begin{cases}
255, & Y(i,j)\geq\mu_G(i,j)-3,\\
0, & Y(i,j)<\mu_G(i,j)-3.
\end{cases}\label{eq:gaussian_bw_threshold}
\end{align}

The important difference from the global methods is that $\mu_G(i,j)$
is not one constant for the entire image. It is a spatially varying
reference calculated from the surrounding grayscale intensities.
\vspace*{\baselineskip}
\subsection{Computational Consideration}

The Gaussian method provides a locally varying reference, but this comes
at an increased computational cost compared with a single global mean.
Conceptually, the local reference has to be evaluated throughout the
image by considering the surrounding pixels and their Gaussian
contributions. The process can therefore be viewed as
\begin{align}
\text{pixel}\;\rightarrow\;\text{neighbourhood}
\;\rightarrow\;\text{Gaussian weights}
\;\rightarrow\;\text{local reference}
\;\rightarrow\;\text{threshold}\label{eq:gaussian_pipeline}
\end{align}

A larger effective Gaussian neighbourhood also involves a broader set
of surrounding pixels. Thus, although the Gaussian reference provides
better spatial adaptation, a direct pixel-by-pixel implementation can be
time-consuming. The \texttt{scipy.ndimage.gaussian\_filter} function used
here performs the filtering efficiently while implementing the same
Gaussian smoothing principle.
